\chapter{Oferta comercial y presupuesto}

El presupuesto se formula como estimacion academica para una oferta preliminar. Incluye hardware, sensores, accesorios, software, ingenieria, instalacion, formacion, mantenimiento, desplazamientos/logistica y margen comercial. Los importes no son una cotizacion real ni sustituyen una consulta a proveedores.

\section{Hipotesis de coste}

\begin{itemize}
  \item Plataforma AMR omnidireccional industrial o equivalente representable en CoppeliaSim: 55.000 \texteuro{} por robot.
  \item Paquete de sensores por robot: LiDAR 2D/3D, RGB-D, RFID/QR, IMU y encoders: 18.500 \texteuro{}.
  \item Accesorios de kitting y portacargas: 6.500 \texteuro{} por robot.
  \item Seguridad industrial adicional: escaner, bumpers y parada de emergencia: 7.500 \texteuro{} por robot.
  \item HMI, tablet y comunicaciones industriales: 2.500 \texteuro{} por robot.
\end{itemize}

\section{Presupuesto por escenarios}

\begin{longtable}{P{0.34\textwidth}r r r}
\toprule
\textbf{Partida presupuestaria} & \textbf{1 robot} & \textbf{7 robots} & \textbf{17 robots}\\
\midrule
\endhead
Plataformas AMR omnidireccionales & 55.000 & 385.000 & 935.000\\
Sensores: LiDAR, RGB-D, RFID/QR, IMU & 18.500 & 129.500 & 314.500\\
Accesorios portacargas/kitting & 6.500 & 45.500 & 110.500\\
Seguridad: escaner, bumpers, E-stop & 7.500 & 52.500 & 127.500\\
HMI, tablets y comunicaciones & 2.500 & 17.500 & 42.500\\
Software de supervision y base de datos & 35.000 & 52.000 & 78.000\\
Ingenieria, simulacion y configuracion & 30.000 & 75.000 & 145.000\\
Instalacion y puesta en marcha & 12.000 & 36.000 & 85.000\\
Formacion de operarios y mantenimiento & 6.000 & 15.000 & 28.000\\
Mantenimiento primer año & 7.200 & 50.400 & 122.400\\
Desplazamientos y logistica & 4.000 & 12.000 & 24.000\\
\midrule
\textbf{Subtotal antes de margen} & \textbf{184.200} & \textbf{870.400} & \textbf{2.012.400}\\
Margen comercial estimado (12\%) & 22.104 & 104.448 & 241.488\\
\midrule
\textbf{Total estimado} & \textbf{206.304} & \textbf{974.848} & \textbf{2.253.888}\\
\bottomrule
\caption{Presupuesto estimado por escenario, en euros.}
\label{tab:presupuesto}
\end{longtable}

\section{Interpretacion comercial}

El escenario de 1 robot se recomienda como piloto controlado y no como solucion definitiva. Su objetivo comercial es reducir incertidumbre tecnica y operativa. El escenario de 7 robots representa una primera expansion industrial con gestor de flota y cobertura de varias estaciones. El escenario de 17 robots se justifica solo si la planta confirma la necesidad de soportar un volumen cercano a 80 HTP al mes o mas, y si los indicadores del piloto y la expansion demuestran utilidad.

La diapositiva extra de presupuesto se incluye en la presentacion del proyecto, con la misma estructura de partidas para que la defensa oral mantenga coherencia con esta oferta comercial.

\section{Coste total de propiedad}

Para evitar una lectura incompleta del presupuesto, se recomienda evaluar el coste total de propiedad a tres años. La estimacion incluye mantenimiento anual, sustitucion preventiva de consumibles, soporte de software y pequeñas adaptaciones de rutas. No se incluyen ampliaciones mayores de integracion MES/ERP ni cambios de layout no previstos.

\begin{longtable}{P{0.36\textwidth}r r r}
\toprule
\textbf{Concepto TCO 3 años} & \textbf{1 robot} & \textbf{7 robots} & \textbf{17 robots}\\
\midrule
\endhead
Inversion inicial estimada & 206.304 & 974.848 & 2.253.888\\
Mantenimiento años 2 y 3 & 14.400 & 100.800 & 244.800\\
Soporte software años 2 y 3 & 12.000 & 28.000 & 42.000\\
Recalibraciones y ajustes de ruta & 5.000 & 18.000 & 36.000\\
\midrule
\textbf{TCO 3 años estimado} & \textbf{237.704} & \textbf{1.121.648} & \textbf{2.576.688}\\
\bottomrule
\end{longtable}

\section{Condiciones comerciales academicas}

Se propone un pago por hitos: 30\% al inicio de ingenieria y simulacion, 40\% tras entrega del piloto instalado, 20\% tras validacion de KPIs y 10\% al cierre documental y formacion. En un contrato real, estos hitos deberian acompañarse de clausulas de aceptacion, responsabilidad, soporte, ciberseguridad y tratamiento de datos.
